\documentclass[9.5pt, a4paper]{article}

% ── PACKAGES ──────────────────────────────────────────────────
\usepackage[utf8]{inputenc}
\usepackage[T1]{fontenc}
\usepackage{geometry}
\geometry{
  top=1.0cm,
  bottom=1.0cm,
  left=1.4cm,
  right=1.4cm
}
\usepackage{hyperref}
\hypersetup{
  colorlinks=true,
  urlcolor=blue,
  linkcolor=blue,
  pdftitle={Identity Axis Framework — One Page Summary},
  pdfauthor={Eric Robert Lawson / OrganismCore}
}
\usepackage{microtype}
\usepackage{parskip}
\usepackage{booktabs}
\usepackage{array}
\usepackage{tabularx}
\usepackage{xcolor}
\usepackage{mdframed}
\usepackage{multicol}
\usepackage{enumitem}
\usepackage{titlesec}
\usepackage{lmodern}

% ── COLOURS ───────────────────────────────────────────────────
\definecolor{headerblue}{RGB}{20,60,140}
\definecolor{statementbox}{RGB}{20,60,140}
\definecolor{rulecolor}{RGB}{20,60,140}

% ── SECTION FORMAT ────────────────────────────────────────────
\titleformat{\section}
  {\small\bfseries\color{headerblue}}
  {}{0em}{}[\vspace{-0.25em}
  \textcolor{rulecolor}{\rule{\linewidth}{0.4pt}}]
\titlespacing{\section}{0pt}{0.45em}{0.2em}

\setlength{\parskip}{0.2em}
\setlength{\parindent}{0pt}
\setlength{\columnsep}{0.9em}
\pagestyle{empty}

% ════════════════════════════════════════��═════════════════════
\begin{document}

% ── HEADER ────────────────────────────────────────────────────
\begin{center}
  {\large\bfseries\color{headerblue}
  THE IDENTITY AXIS FRAMEWORK}\\[0.15em]
  {\small\color{headerblue}
  A Geometric Model of Cancer — One-Page Summary}\\[0.1em]
  {\footnotesize
  Eric Robert Lawson\ $\cdot$\
  OrganismCore (Independent Research)\ $\cdot$\
  ORCID:
  \href{https://orcid.org/0009-0002-0414-6544}
  {0009-0002-0414-6544}\ $\cdot$\
  \href{mailto:OrganismCore@proton.me}
  {OrganismCore@proton.me}\ $\cdot$\
  2026-03-07}
\end{center}

\vspace{0.15em}

% ── THE CLAIM ─────────────────────────────────────────────────
\begin{mdframed}[
  backgroundcolor=statementbox,
  linecolor=statementbox,
  linewidth=0pt,
  innertopmargin=0.35em,
  innerbottommargin=0.35em,
  innerleftmargin=0.7em,
  innerrightmargin=0.7em
]
\begin{center}
\footnotesize\bfseries\color{white}
Cancer is a cell displaced from its normal attractor in the Waddington
epigenetic landscape and trapped in a false one.
The geometry of that displacement is measurable.
The drug target is derivable from the cell of origin — before any
literature is consulted.
Demonstrated across \textbf{41 cancer types} with
\textbf{zero structural contradictions}.
\end{center}
\end{mdframed}

\vspace{0.3em}

% ── TWO COLUMNS ───────────────────────────────────────────────
\begin{multicols}{2}

% ── LEFT COLUMN ───────────────────────────────────────────────

\section{The Theorem}

Every cancer has a measurable axis between:

\smallskip
{\small
\begin{tabular}{p{0.55cm}p{5.5cm}}
\textbf{A} &
\textbf{Identity Anchor} — the gene defining
what the cell of origin \textit{is}. \\[0.2em]
\textbf{H} &
\textbf{Convergence Hub} — the gene defining
what the cancer has \textit{committed to}. \\
\end{tabular}
}

\smallskip
{\small
$R = A/H$ is a continuous, IHC-measurable
coordinate in the Waddington landscape.
$R$ predicts overall survival.
\textbf{H is the drug target.}
}

{\footnotesize
DOI:\
\href{https://doi.org/10.5281/zenodo.18898788}
{\texttt{10.5281/zenodo.18898788}}
}

\section{The Evidence}

{\small
\begin{tabular}{p{4.0cm}p{1.8cm}}
Cancer types analysed & \textbf{41} \\
Structural contradictions & \textbf{0} \\
FDA drugs confirming geometry-H & \textbf{5} \\
Phase 1--3 drugs confirming H & \textbf{14} \\
BRCA datasets validated & \textbf{7} \\
BRCA patients validated & \textbf{$\sim$7,500} \\
Novel predictions pre-registered & \textbf{14} \\
\end{tabular}
}

\section{The Generative Proof}

{\small
Five geometric coordinates were
\textit{constructed} without starting from
any known cancer. Five cancers were derived
and confirmed with exactly the predicted
properties. \textbf{The cancer was the
output. The geometry was the input.}
This eliminates the pattern-matching
objection. Cancer is calculable.
}

{\footnotesize
DOI:\
\href{https://doi.org/10.5281/zenodo.18899978}
{\texttt{10.5281/zenodo.18899978}}
}

\section{The Four Root Documents}

{\footnotesize
\begin{tabular}{p{3.6cm}p{2.4cm}}
\toprule
\textbf{Document} & \textbf{DOI} \\
\midrule
Identity Axis Theorem &
\href{https://doi.org/10.5281/zenodo.18898788}
{\texttt{...18898788}} \\[0.15em]
Cancer Is Calculable &
\href{https://doi.org/10.5281/zenodo.18899250}
{\texttt{...18899250}} \\[0.15em]
Tazemetostat Indication &
\href{https://doi.org/10.5281/zenodo.18899530}
{\texttt{...18899530}} \\[0.15em]
Cancer Via Construction &
\href{https://doi.org/10.5281/zenodo.18899978}
{\texttt{...18899978}} \\
\bottomrule
\end{tabular}

\smallskip
Repository:\
\href{https://github.com/Eric-Robert-Lawson/attractor-oncology}
{\texttt{github.com/Eric-Robert-Lawson/}
\texttt{attractor-oncology}}
}

\columnbreak

% ── RIGHT COLUMN ���─────────────────────────────────────────────

\section{Selected Cancer Axes}

{\small
\begin{tabular}{p{1.6cm}p{1.9cm}p{1.7cm}}
\toprule
\textbf{Cancer} &
\textbf{A} &
\textbf{H} \\
\midrule
Breast & FOXA1 & EZH2 \\
ccRCC & GOT1 & RUNX1 \\
AML & CEBPA & HOXA9 \\
SCLC-Y & ASCL1 & REST \\
DIPG$^\dagger$ & OLIG2 & H3K27M \\
Uveal mel.$^\ddagger$ & MITF & BAP1 \\
Ewing & RUNX2/SOX9 & EZH2 \\
Thymoma/TC & FOXN1 & EZH2 \\
Merkel cell & ATOH1 & EZH2 \\
Angiosarcoma & ERG & EZH2 \\
\midrule
\multicolumn{3}{p{5.4cm}}{
\footnotesize
H = EZH2 in 26/41 (72\%) of cancers.
In the remaining 28\%, the framework
identifies the specific Polycomb arm
operative in that lineage.}\\[0.15em]
\multicolumn{3}{p{5.4cm}}{
\footnotesize
$^\dagger$Mirror case: EZH2 paralysed
by H3K27M. Tazemetostat
contraindicated.}\\[0.1em]
\multicolumn{3}{p{5.4cm}}{
\footnotesize
$^\ddagger$PRC1 arm operative.
BAP1 loss = attractor deepening.}\\
\bottomrule
\end{tabular}
}

\section{The FOXA1/EZH2 Clinical
Translation — One Step Remaining}

{\small
\begin{tabular}{p{2.7cm}p{2.7cm}}
\textbf{Exists} & \textbf{Needed} \\
\midrule
7 validated datasets & IHC cut-points \\
$\sim$7,500 patients & 300--400 FFPE cases \\
Survival HR confirmed & PAM50 ground truth \\
Protocol published & One pathology partner \\
Cost: \$50--\$100/pt & H-score scoring \\
\end{tabular}
}

\smallskip
{\small
The calibration study translates a
computationally validated classifier
into a \$50 breast cancer subtype tool
deployable in any pathology laboratory
globally — on standard FFPE tissue,
no RNA, no proprietary platform,
no capital equipment.

\textbf{1.2--1.5 million patients per year
currently receive no molecular subtype
information. This test reaches all of them.}

The first institution to run the
calibration study is first co-author
on the publication.
}

\end{multicols}

\vspace{0.1em}
\noindent\textcolor{rulecolor}
{\rule{\linewidth}{0.4pt}}

% ── FOOTER ────────────────────────────────────────────────────
\begin{center}
\footnotesize
\textit{The geometry came first.\ \
The data confirmed it.\ \
The drugs follow from the geometry.\ \
The patients are the reason for all of it.}\\[0.2em]
\href{https://github.com/Eric-Robert-Lawson/attractor-oncology}
{\texttt{github.com/Eric-Robert-Lawson/attractor-oncology}}
\ $\cdot$\
\href{https://doi.org/10.5281/zenodo.18898788}
{\texttt{doi.org/10.5281/zenodo.18898788}}
\ $\cdot$\
\href{https://orcid.org/0009-0002-0414-6544}
{\texttt{ORCID: 0009-0002-0414-6544}}
\end{center}

\end{document}
